\documentclass[11pt,aps,preprint,nofootinbib,superscriptaddress,tightenlines]{revtex4-2}
\usepackage{amsmath,amssymb,bm,mathtools,graphicx,dsfont}
\usepackage{xcolor}
\definecolor{darkblue}{RGB}{0,0,160}
\usepackage[colorlinks=true,linkcolor=darkblue,citecolor=darkblue,urlcolor=darkblue]{hyperref}

\def\st{\begin{equation}}
\def\stp{\end{equation}}
\def\Eq#1{eq.~(\ref{#1})}
\def\Sect#1{Sect.~\ref{#1}}
\newcommand{\one}{\mathds{1}}
\def\tr{\operatorname{tr}}
\let\oldparagraph\paragraph
\renewcommand{\paragraph}{\par\bigskip\oldparagraph}

\begin{document}

\title{One comoving kernel for the ideal and the viscous step}
\author{D.~Teaney}
\affiliation{Department of Physics and Astronomy, Stony Brook University,
Stony Brook, New York 11794, USA}
\date{\today}

\begin{abstract}
\texttt{commonkernel.tex} makes the two steps share one kernel by moving the viscous step to the lab.
Here we go the other way: the viscous step is kept as it is, and the ideal step measures its density with
the same comoving kernel.  The density depends on the positions and on the particle's own velocity only
through rotation invariants, so the Hamiltonian is invariant under translations and rotations.  Taking all its
derivatives, including those through the velocity, gives equations of motion that conserve energy, momentum
and angular momentum.  The particles then drift at the density-frame velocity corrected by the momentum of the
kernel's gradient stress, as the density frame requires.  This note defines the step and lists what must be
checked.  Nothing here has been run.
\end{abstract}

\maketitle

%=============================================================================
\section{Notation}
\label{sec:notation}

As in \texttt{commonkernel.tex}: $\bm r_{ab}=\bm x_b-\bm x_a$, $V_a$ the particle volume, $s_a$ its entropy,
$\bm m_a$ its momentum, and $e_a(s_a,\bm m_a,V_a)$ its energy.  The velocity, pressure and comoving metric are
\st
\bm v_a=\frac{\partial e_a}{\partial\bm m_a}\Big|_{s,V},\qquad
p_a=-\frac{\partial e_a}{\partial V_a}\Big|_{s,\bm m},\qquad
\ell_a=\one-\bm v_a\bm v_a^T,\qquad \ell_a^{-1}=\one+\gamma_a^2\bm v_a\bm v_a^T .
\stp
The comoving distance of the bond $ab$, seen from $a$, is
\st
\rho_{ab}^2=\bm r_{ab}\cdot\ell_a^{-1}\bm r_{ab}=r_{ab}^2+\gamma_a^2(\bm v_a\cdot\bm r_{ab})^2 .
\stp
$W(\rho;h)$ is the SPH kernel and
\st
\psi(\rho)=-\frac{W'(\rho)}{\rho}
\stp
its bond weight.

%=============================================================================
\section{The density}
\label{sec:density}

Particle $a$ measures its number density with its own comoving kernel,
\st
\boxed{\;
n_a=N_a\big(\{\bm r_{ab}\},\bm v_a\big)\equiv\gamma_a\sum_bW(\rho_{ab}),\qquad
V_a=\frac1{n_a},\qquad H=\sum_ae_a(s_a,\bm m_a,V_a) .\;}
\label{H}
\stp
The sum includes $b=a$.  The factor $\gamma_a$ normalizes the kernel,
\st
\int d^2r\;\gamma_aW(\rho)=1 ,
\stp
so for smooth fields $n_a$ is the \emph{lab} density, as the density-frame energy requires.  Since $\bm v_a$ depends on $V_a$,
\Eq{H} is a scalar equation for $n_a$ per particle, solved by iteration (one or two passes from the previous step).

There is no moment correction in the density: \Sect{sec:closure} shows that it would not improve the gradient.  The correction
$\Lambda_a$ stays in the viscous step.

\paragraph{Invariances.}  Under a translation of all positions nothing changes.  Under a rotation
\st
\bm x_a\to R\bm x_a,\qquad \bm m_a\to R\bm m_a ,
\stp
$e_a$ depends on $|\bm m_a|$, so $\bm v_a\to R\bm v_a$ and $\ell_a\to R\ell_aR^T$, and $\rho_{ab}$ is invariant.  $H$ is invariant under
both, and its equations of motion conserve
\st
H,\qquad \bm P=\sum_a\bm m_a,\qquad L=\sum_a\bm x_a\times\bm m_a .
\stp

%=============================================================================
\section{The equations of motion}
\label{sec:eom}

\paragraph{The gradient.}  Differentiating \Eq{H},
\st
\bm D_{ab}\equiv-\frac{\partial N_a}{\partial\bm r_{ab}}=\gamma_a\,\psi(\rho_{ab})\,\ell_a^{-1}\bm r_{ab} .
\label{D}
\stp
At rest this is the present kernel gradient.  The bond weight enters the ideal step only here.

\paragraph{The velocity derivative.}  From
\st
\frac{\partial\gamma_a}{\partial v^k}=\gamma_a^3\,v_k,\qquad
\frac{\partial\rho_{ab}^2}{\partial v^k}=2\gamma_a^2\,\big(\ell_a^{-1}\bm r_{ab}\big)_k\,(\bm r_{ab}\cdot\bm v_a) ,
\stp
the derivative of the density with respect to the velocity is
\st
\frac{\partial n_a}{\partial v^k}\Big|_{\bm r}=-\gamma_a^2\Big[\sum_bD_{ab\,k}\,(\bm r_{ab}\cdot\bm v_a)-n_a\,v_{a\,k}\Big] .
\label{dndv}
\stp
The bracket is the closure defect of the gradient, \Eq{closure}, acting on the velocity.  For a uniform density it vanishes, because
the kernel is normalized at every velocity.

\paragraph{The momentum of the kernel stress.}  With
\st
P_a\equiv f_a\,p_aV_a^2
\stp
($f_a$ from \Eq{fa} below, $f_a=1$ at rest), define the covector
\st
\boxed{\;
\Theta^0_{a\,k}\equiv-\frac{P_a}{V_a}\,\frac{\partial n_a}{\partial v^k}
=\frac{P_a}{V_a}\,\gamma_a^2\Big[\sum_bD_{ab\,k}\,(\bm r_{ab}\cdot\bm v_a)-n_a\,v_{a\,k}\Big] .\;}
\label{Theta0}
\stp
It has the index of a momentum density.  \Sect{sec:densityframe} shows that in the continuum it is
\st
\Theta^0_k=\Theta_{kl}\,v^l,\qquad \Theta_{kl}=\sigma^2\,\frac{p}{n}\,\partial_k\partial_ln ,
\stp
the lab momentum density of the gradient stress of the kernel, which is purely spatial in the kernel's rest frame.

\paragraph{The variation of $n_a$.}  The velocity depends on the momentum and on the volume,
\st
\delta\bm v_a=\frac{\partial\bm v_a}{\partial\bm m_a}\,\delta\bm m_a+\frac{\partial\bm v_a}{\partial V_a}\,\delta V_a,\qquad
\delta V_a=-V_a^2\,\delta n_a ,
\stp
so
\st
\delta n_a=f_a\Big[\sum_b\bm D_{ab}\cdot(\delta\bm x_a-\delta\bm x_b)+\frac{\partial n_a}{\partial\bm v_a}\cdot\frac{\partial\bm v_a}{\partial\bm m_a}\,\delta\bm m_a\Big],\qquad
f_a=\Big[1+V_a^2\,\frac{\partial n_a}{\partial\bm v_a}\cdot\frac{\partial\bm v_a}{\partial V_a}\Big]^{-1} .
\label{fa}
\stp

\paragraph{Hamilton's equations.}  With
\st
\delta e_a=\bm v_a\cdot\delta\bm m_a+p_aV_a^2\,\delta n_a
\stp
at fixed $s_a$, and
\st
\frac{\partial\bm v_a}{\partial\bm m_a}=\frac1{V_a}\,\frac{\partial\bm v}{\partial\bm M},\qquad \bm M=\frac{\bm m}{V},
\stp
\st
\boxed{\;
\dot{\bm x}_a=\bm v_a-\frac{\partial\bm v}{\partial\bm M}\,\bm\Theta^0_a,\qquad
\dot{\bm m}_a=-\sum_{b\neq a}\Big[P_a\,\bm D_{ab}-P_b\,\bm D_{ba}\Big],\qquad
\dot s_a=0 .\;}
\label{eom}
\stp
The force has the present pair form, with $\bm D_{ab}\neq-\bm D_{ba}$ when the two metrics differ; momentum is conserved pair by pair.
The particle drifts at the density-frame velocity, minus the inverse susceptibility acting on the momentum of the kernel stress.

\paragraph{Angular momentum.}  In 2D write
\st
\bm a\times\bm b=\epsilon_{ij}a^ib^j,\qquad (\epsilon\bm a)_i=\epsilon_{ij}a^j,\qquad \bm a\cdot\epsilon\bm b=\bm a\times\bm b .
\stp
Two identities follow from an infinitesimal rotation $R=\one+\theta\epsilon$.
\begin{itemize}
\item $N_a$ is invariant when all bonds and the velocity are rotated.  At first order in $\theta$,
  \st
  \sum_b\bm D_{ab}\times\bm r_{ab}=\frac{\partial n_a}{\partial\bm v_a}\times\bm v_a .
  \label{wardN}
  \stp
\item $e_a$ depends on $|\bm m_a|$ only, so the velocity is along the momentum and rotates with it:
  \st
  \bm v_a\times\bm m_a=0,\qquad \frac{\partial\bm v_a}{\partial\bm m_a}\,\epsilon\bm m_a=\epsilon\bm v_a .
  \label{wardv}
  \stp
\end{itemize}
The rate of change of the angular momentum is
\st
\dot L=\sum_a\big(\dot{\bm x}_a\times\bm m_a+\bm x_a\times\dot{\bm m}_a\big) .
\stp
In the force term, relabel $a\leftrightarrow b$ in the $P_b$ part of \Eq{eom} and use \Eq{wardN}:
\st
\sum_a\bm x_a\times\dot{\bm m}_a=\sum_aP_a\sum_b\bm r_{ab}\times\bm D_{ab}
=-\sum_aP_a\,\frac{\partial n_a}{\partial\bm v_a}\times\bm v_a
=\sum_aV_a\,\bm\Theta^0_a\times\bm v_a .
\stp
In the drift term use \Eq{wardv} and the symmetry of the velocity derivative:
\st
\sum_a\dot{\bm x}_a\times\bm m_a=-\sum_aV_a\Big(\frac{\partial\bm v_a}{\partial\bm m_a}\bm\Theta^0_a\Big)\cdot\epsilon\bm m_a
=-\sum_aV_a\,\bm\Theta^0_a\cdot\epsilon\bm v_a
=-\sum_aV_a\,\bm\Theta^0_a\times\bm v_a .
\stp
The two cancel star by star, and
\st
\dot L=0 .
\stp
The non-central pair forces put the torque
\st
V_a\,\bm\Theta^0_a\times\bm v_a
\stp
on the star of particle $a$, and the drift correction carries exactly the opposite angular momentum.  With the metric frozen, the
drift correction is absent but the torque is not.

In a periodic box there is no rotation symmetry and $L$ is not conserved globally; what survives is the torque balance star by star,
and the clean test is an isolated rotating patch.

%=============================================================================
\section{The closure}
\label{sec:closure}

Deform the whole star of $a$ linearly, at fixed velocity,
\st
\bm r_{ab}\to A\,\bm r_{ab} .
\stp
Since $\bm D_{ab}$ is minus the derivative of $N_a$ with respect to $\bm r_{ab}$, its first moment is exactly
\st
\sum_b\bm D_{ab}\,\bm r_{ab}^T=-\frac{\partial N_a\big(\{A\bm r_{ab}\},\bm v_a\big)}{\partial A}\Big|_{A=\one} .
\label{closure}
\stp
For smooth fields the density scales as $1/\det A$, and with
\st
\int d^2r\;\psi(\rho)\,\bm r\bm r^T=\frac{\ell_a}{\gamma_a}
\stp
the first moment is
\st
\sum_b\bm D_{ab}\,\bm r_{ab}^T\simeq n_a\,\one .
\stp
On a disordered star this is not exact.  The moment correction of the viscous step, put into the density, does not fix it: the
terms from the correction's own dependence on the bonds are as large as the error it removes.  On a jittered lattice (Wendland C2,
$h/a=2$, jitter $0.15a$), with the correction re-solved in each finite difference, the closure divided by $n_a$ is
\begin{center}
\begin{tabular}{lcc}
velocity & this density & with $\Lambda_a$ in the density\\ \hline
$0$ & diag$(0.947,\,1.027)$ & diag$(0.959,\,1.032)$\\
$(0.5,\,0.2)$ & diag$(0.920,\,1.024)$ & diag$(0.927,\,1.029)$\\
$(0.8,\,0)$ & diag$(0.910,\,1.050)$ & diag$(0.885,\,1.052)$
\end{tabular}
\end{center}
with off-diagonal entries below $0.05$.  This is one star; over five jittered stars the diagonal entries range from $0.80$ to $1.11$, at
every velocity.  In the viscous step the closure
\st
\sum_b\bm A_{ab}\,\bm r_{ab}^T=V_a\,\one
\stp
is exact because the faces are built directly.  A Hamiltonian gradient is the derivative of a density, so its closure is a property of
that density; exact closure would need a density that scales exactly as $1/\det A$ under the deformation.  The closure error of
\Eq{D} is that of ordinary SPH, now comoving.  The present ideal step has no sampling correction either: its density and force use
the plain kernel and its gradient (\texttt{smooth\_entropy\_energy!}, \texttt{hamiltonian\_force\_energy!}), and at rest \Eq{D} is that gradient.
So nothing is lost relative to the present step.

%=============================================================================
\section{Integration}
\label{sec:integration}

$H$ is already non-separable ($e_a$ depends on the momentum and on the volume), and the present drift--kick--drift
(\texttt{leapfrog\_step!}) applies unchanged with the rates \Eq{eom}.  It conserves $\bm P$ exactly.  The present scheme conserves $L$
exactly in each substep (central forces, drift along the momentum); here the drift and the kick each change $L$.  With
\st
X(\bm x,\bm m)\equiv\sum_aV_a\,\bm\Theta^0_a\times\bm v_a ,
\stp
a drift of $\delta t$ at fixed momenta changes $L$ by exactly $-\delta t\,X$, and a kick by exactly $+\delta t\,X$, so one step changes it by
\st
\Delta L=\Delta t\Big[-\tfrac12X(\bm x^n,\bm m^n)+X(\bm x^{n+1/2},\bm m^n)-\tfrac12X(\bm x^{n+1/2},\bm m^{n+1})\Big] .
\stp
If the kick were evaluated at $\bm m^{n+1/2}$ this would be a second difference, of order $\Delta t^3$.  The present kick uses $\bm m^n$, and
since the pressure depends on the momentum the bracket is of order $\Delta t$: $\Delta L$ is of order $\Delta t^2$ per step.  The same holds for
$H$ in the present step.  To be measured; an implicit midpoint step would conserve $L$ exactly.

%=============================================================================
\section{The continuum limit and the density frame}
\label{sec:densityframe}

\paragraph{The smoothed density.}  For smooth fields, expanding the density to second order about $\bm x_a$,
\st
n_a\simeq n+\tfrac12\sigma^2\,\ell_a^{ij}\,\partial_i\partial_jn,\qquad
\sigma^2=\tfrac12\int d^2u\;W(u)\,u^2 ,
\stp
since the second moment of the normalized comoving kernel is $\sigma^2\ell_a$.  For the Wendland kernels
\st
\frac{\sigma^2}{h^2}=\frac5{18}\ (\text{C2}),\qquad \frac7{33}\ (\text{C4}),\qquad \frac6{35}\ (\text{C6}).
\stp
The Hamiltonian therefore gains a reversible gradient term
\st
H\simeq H_0+\sum_aV_a\,\tfrac12\,\ell_a^{ij}\,\Theta_{ij},\qquad
\Theta_{ij}\equiv\sigma^2\,\frac{p}{n}\,\partial_i\partial_jn ,
\stp
a stress of order $h^2\nabla\nabla$ (Korteweg-like) that depends on the momentum through $\ell_a$.

\paragraph{Indices.}  The indices sit where they are born: $\Theta_{ij}$ is lower (two derivatives), $\ell^{ij}$ is upper,
\st
\ell^{ij}=e^i_\mu e^j_\nu\,\eta^{\mu\nu}=\delta^{ij}-v^iv^j
\stp
(the comoving inverse metric of \texttt{theory.tex}), $v^i$ is upper, and the momentum density $M_k$ is lower.  With
\st
\frac{\partial\ell^{ij}}{\partial v^k}=-\delta^i_kv^j-v^i\delta^j_k ,
\stp
\Eq{dndv} and \Eq{Theta0} become
\st
\frac{\partial n_a}{\partial v^k}=-\sigma^2\,\partial_k\partial_jn\;v^j,\qquad
\Theta^0_k=\Theta_{kj}\,v^j ,
\stp
and the drift is
\st
\dot x^i-v^i=-\frac{\partial v^i}{\partial M_k}\,\Theta_{kj}\,v^j .
\label{dx}
\stp
No metric is used: the free index is raised by the inverse susceptibility.

\paragraph{The inverse susceptibility.}  The derivative is at fixed $s_a$ and $V_a$, i.e.\ at fixed entropy density $S$.  With
\st
v^i=\frac{M^i}{E+\mathcal P},\qquad \frac{\partial E}{\partial M_k}\Big|_S=v^k ,
\stp
\st
\frac{\partial v^i}{\partial M_k}\Big|_S=\frac{G^{ik}}{E+\mathcal P},\qquad
G^{ik}=\underbrace{\delta^{ik}-v^iv^k}_{h^{ik}}\;-\;v^i\,\frac{\partial\mathcal P}{\partial M_k}\Big|_S
=h^{ik}+\frac{c_s^2}{\gamma^2}\,\frac{v^iv^k}{1-c_s^2v^2},\qquad
E+\mathcal P=\gamma^2TS .
\label{G}
\stp
The comoving metric comes from the momentum and the energy; the second term is the pressure response,
\st
\frac{\partial\mathcal P}{\partial M_k}\Big|_S=\frac{\partial\mathcal P}{\partial M_k}\Big|_E+v^k\frac{\partial\mathcal P}{\partial E}\Big|_{\bm M}
=-\frac{2c_s^2\,v^k}{1-c_s^2v^2}+\frac{c_s^2(1+v^2)\,v^k}{1-c_s^2v^2}
=-\frac{c_s^2}{\gamma^2}\,\frac{v^k}{1-c_s^2v^2},
\stp
where the two derivatives are the $(0k)$ (times two) and $(00)$ components of the pressure derivatives in \texttt{theory.tex},
eq.~(pressure\_derivs).  (Checked against a numerical Jacobian at fixed $S$.)

The sign is that of a boost at fixed entropy.  With $s$, $T$ and $w=Ts$ the rest-frame entropy density, temperature and enthalpy, fixed
$S=\gamma s$ means
\st
\frac{ds}{s}=-\frac{d\gamma}{\gamma}=-\gamma^2\,\bm v\cdot d\bm v ,
\stp
so the fluid cools and its pressure drops:
\st
d\mathcal P=s\,dT=c_s^2\,w\,\frac{ds}{s}=-c_s^2\,w\gamma^2\,\bm v\cdot d\bm v=-c_s^2\,v_i\,G^{ik}\,dM_k ,\qquad
v_iG^{ik}=\frac{v^k}{\gamma^2(1-c_s^2v^2)} .
\stp
Of the two pieces, the derivative at fixed $E$ is negative (more momentum leaves less internal energy) and the second is positive
(keeping $S$ fixed costs energy); the first wins by a factor $2/(1+v^2)$.

\paragraph{The momentum of the kernel stress.}  $\Theta$ is a stress defined in the rest frame of the kernel, so it has no temporal
components there,
\st
\Theta_{\mu\nu}u^\nu=0 .
\stp
In the lab, with $u^\nu=\gamma(1,\bm v)$, this gives it the momentum density
\st
\Theta_{k0}=-\Theta_{kl}\,v^l ,\qquad \Theta^0{}_k=\Theta_{kl}\,v^l ,
\stp
which is $\bm\Theta^0_a$ of \Eq{Theta0}.  The density frame allows no correction to $T^{0k}$, so this momentum is absorbed by a frame
change (\texttt{theory.tex}),
\st
\delta\beta_\rho=-\chi^{-1}_{\rho\nu}\,\underline\Pi^{0\nu} ,
\stp
and the density-frame velocity, read off the momentum, differs from the velocity of the rest frame of the stress by \Eq{dx}.  So the
particles, which carry the entropy, move with the rest frame in which the stress is purely spatial (the Landau velocity of $\Theta$),
and the velocity of the equation of state is the density-frame velocity.  (Up to $\Theta^{00}=\Theta^{jk}v_jv_k$, which shifts $E$ and is of
second order in $v$; to be checked with the momentum flux.)

\paragraph{The entropy flux.}  In the density frame (\texttt{vischydro/theory.tex}, Sect.~2) the densities $T^{00}$, $T^{0i}$ are ideal
and every correction sits in the spatial fluxes: a stress $\Pi^{ij}$ in $T^{ij}$ comes with the entropy flux
\st
J_S^j=Sv^j-\beta_i\Pi^{ij},\qquad \beta_i=\delta_{ik}\frac{v^k}{T} .
\stp
In the particle scheme the entropy is carried at the drift velocity, so with \Eq{dx} and \Eq{G}
\st
S\big(\dot x^i-v^i\big)=-\frac1T\,\frac{G^{ik}}{\gamma^2}\,\Theta_{kj}\,v^j ,
\stp
which is $-\beta_j\Pi^{ji}$ with
\st
\Pi^{ij}v_j=\frac{G^{ik}}{\gamma^2}\,\Theta_{kj}\,v^j .
\stp
At small $v$ this is $\Pi^{ij}\simeq\delta^{ik}\delta^{jl}\Theta_{kl}$.  At finite $v$ the entropy flux fixes only the contraction with $v$; the full
stress must be read off the momentum flux of $H$.  This is the mechanism of the frame-change matrix $\kappa^{ij}_{\mu\nu}$ of
\texttt{theory.tex}: a rest-frame stress written in lab densities acquires the advective factor
\st
e^i_\mu=\delta^i_\mu-v^i\delta^0_\mu
\stp
and the pressure response.  Here the stress is reversible and of order $h^2\nabla\nabla$, not viscous, and the Hamiltonian supplies its
density-frame form through the drift.

\paragraph{To check analytically:} the momentum flux of $H$ to order $h^2$, which gives the full stress; whether it equals
$\kappa^{ij}_{\mu\nu}$ applied to the rest-frame stress built from $\Theta$; and that its contraction with $v$ reproduces the entropy flux
above.

%=============================================================================
\section{The viscous and stochastic steps}
\label{sec:viscous}

Unchanged: the faces
\st
\bm A_{ab}=V_a\,w_{ab}\,\ell_a^{-1}\bm r_{ab}
\stp
with the $\Lambda_a$-corrected weights of \texttt{stochastic.tex}, the fit, the stress move, the residual move and the vote.  For Wendland
C2 with $h_\psi=h$ the bond weight is the present one; for C4 and C6 it becomes the bond weight of those kernels.  Both steps then use the
same $\psi(\rho)$ and the same metric.  Up to normalization and the correction, the face is the gradient of the ideal step,
\st
\bm A_{ab}\simeq V_aV_b\,\bm D_{ab} .
\stp

%=============================================================================
\section{The cost: fewer neighbours along the flow}
\label{sec:cost}

The ellipse $\rho<2h$ has area $\pi(2h)^2/\gamma_a$.  At a fixed lab spacing it holds $1/\gamma_a$ of the neighbours at rest, and its width
along the flow is $2h/\gamma_a$: an effective resolution
\st
\Big(\frac ha\Big)_{\rm along}=\frac{h/a}{\gamma}
\stp
($1.2$ at $h/a=2$, $v=0.8$).  The viscous step already pays this; the ideal step would pay it too.

%=============================================================================
\section{The fixed lattice}
\label{sec:grid}

On the lattice of \texttt{gridstep.tex} the kick and the drift come from the same $H$, so the drift displacement is the drift velocity
of \Eq{eom} times the time step.  The momentum of the kernel stress is nonzero where the velocity or the density varies from site to
site.  The return is unchanged.

%=============================================================================
\section{Implementation plan}
\label{sec:plan}

\paragraph{Scope.}  A new ideal step, selected by \texttt{SPH\_IDEAL=comoving} (\texttt{SimParams.ideal = :comoving}), for the number-volume
energy Hamiltonian in 2D (\texttt{hamiltonian = :energy}, \texttt{volume = :number}), as the grid option is restricted now.  The default
\texttt{:particle} is untouched and stays bit-identical.  The viscous and stochastic steps are not edited.

\paragraph{The particle record.}  \texttt{Particle} keeps its fields.  \texttt{dotx} and \texttt{Teff} are left from the two-level
Hamiltonian, where they were the smoothed $\partial H/\partial\bm m_a$ and $\partial H/\partial s_a$; in the number-volume Hamiltonian
\texttt{update\_velocities\_energy!} sets them to $\bm v_a$ and $T_a$.  The viscous and stochastic steps read \texttt{dotx} as the
velocity of the comoving metric and form $\beta_a=\texttt{dotx}/\texttt{Teff}$; both must stay the density-frame $\bm v_a$ and $T_a$.  So
the drift velocity of \Eq{eom} is computed and used inside the ideal step only, in a scratch array, and never written to
\texttt{dotx}.

\paragraph{The density, shared.}  The viscous and stochastic steps call \texttt{smooth\_to\_particles!}, which has no equation of state.
For \texttt{:comoving} it evaluates \Eq{H} in one pass with the velocity stored from the last solve,
\st
n_a=\gamma_a\sum_bW(\rho_{ab}),\qquad \rho_{ab}^2=r_{ab}^2+\gamma_a^2(\bm v_a\cdot\bm r_{ab})^2,\qquad \bm v_a=\texttt{v\_over\_S}\cdot\texttt{S\_smooth},
\stp
and sets \texttt{S\_smooth} $=s_an_a$, \texttt{M\_smooth} $=\bm m_an_a$ as now.  The ideal step ends on a converged state, so at fixed positions this
pass returns the same density, and the dissipative steps see the volumes of the Hamiltonian.

\paragraph{The fixed point.}  In the ideal step the density and the velocity are made consistent per particle,
\st
n_a^{(k+1)}=N_a\Big(\{\bm r_{ab}\},\ \bm v\big(s_a,\bm m_a,1/n_a^{(k)}\big)\Big) ,
\stp
alternating the pass above with \texttt{update\_thermodynamics!} until
\st
\max_a\big|n_a^{(k+1)}-n_a^{(k)}\big|/n_a<10^{-13} .
\stp
The contraction factor is $1/f_a-1$, of order $\Theta^0$, so two or three passes are expected; the count is reported.

\paragraph{The derivatives of the equation of state.}  The drift and $f_a$ need, at fixed $s_a$,
\st
\frac{\partial\bm v}{\partial\bm M}\Big|_S=\frac{G}{E+\mathcal P}\quad\text{(\Eq{G})},\qquad
\frac{\partial\bm v_a}{\partial V_a}\Big|_{s,\bm m}=-\frac1{V_a}\Big[\frac{\partial\bm v}{\partial\bm M}\Big|_S\bm M+\frac{\partial\bm v}{\partial S}\Big|_{\bm M}S\Big] .
\stp
The first exists (\texttt{inverse\_susceptibility}, up to the factor $T$); the second is new.  Both are written analytically for the conformal
equation of state and tested against finite differences of \texttt{solve\_eos!}.

\paragraph{The sums.}  One loop over the neighbours per particle gives, at the converged velocity,
\st
n_a,\qquad
\bm Q_a\equiv\sum_b\bm D_{ab}\,(\bm r_{ab}\cdot\bm v_a),\qquad
\frac{\partial n_a}{\partial\bm v_a}=-\gamma_a^2\big(\bm Q_a-n_a\bm v_a\big),
\stp
and from them
\st
f_a=\Big[1+V_a^2\,\frac{\partial n_a}{\partial\bm v_a}\cdot\frac{\partial\bm v_a}{\partial V_a}\Big]^{-1},\qquad
P_a=f_a\,p_aV_a^2,\qquad
\bm\Theta^0_a=-\frac{P_a}{V_a}\,\frac{\partial n_a}{\partial\bm v_a},\qquad
\dot{\bm x}_a=\bm v_a-\frac{\partial\bm v}{\partial\bm M}\,\bm\Theta^0_a .
\stp
A second loop gives the force,
\st
\dot{\bm m}_a=-\sum_b\big[P_a\,\bm D_{ab}-P_b\,\bm D_{ba}\big],\qquad
\bm D_{ab}=\gamma_a\,\psi(\rho_{ab})\,\ell_a^{-1}\bm r_{ab},\qquad
\bm D_{ba}=-\gamma_b\,\psi(\rho_{ba})\,\ell_b^{-1}\bm r_{ab} ,
\stp
with $P_b$ from the first loop.

\paragraph{The integrator.}  The drift--kick--drift of \texttt{leapfrog\_step!}, with \texttt{leapfrog\_step!} dispatching to the new step
when \texttt{params.ideal === :comoving}, so that every caller (\texttt{generic\_step!}, \texttt{stability.jl}, \texttt{equilibrium.jl})
gets it without edits.

\paragraph{The hand-off to the dissipative steps.}  This is the one place where ``the dissipative steps are unchanged'' is not exact.
The stochastic step holds $e_a$ and $V_a$ fixed and changes $\bm m_a$.  With the comoving kernel the volume depends on the momentum, so when the
next step recomputes it at the new velocity the Hamiltonian changes by
\st
\delta H\simeq\sum_ap_aV_a^2\,\frac{\partial n_a}{\partial\bm v_a}\cdot\delta\bm v_a=-\sum_aV_a\,\bm\Theta^0_a\cdot\delta\bm v_a ,
\stp
of random sign and of order $\Theta^0$.  It is the work of the kernel stress.  With
\st
\delta\bm v_a=\frac1{V_a}\frac{\partial\bm v}{\partial\bm M}\,\delta\bm m_a ,
\stp
it reads
\st
\delta H\simeq\sum_a\big(\dot{\bm x}_a-\bm v_a\big)\cdot\delta\bm m_a :
\stp
the dissipative steps move momentum as if each particle moved at the density-frame velocity, while the velocity conjugate to the
momentum in $H$ is the drift velocity.  This is the GENERIC degeneracy condition, which needs $\partial H/\partial\bm m_a$ in the
dissipative bracket.

\paragraph{The fix in the accept/reject.}  The volume of particle $a$ depends on its own momentum only, with
\st
p_a\,\frac{\partial V_a}{\partial\bm m_a}=-p_aV_a^2\,f_a\,\frac{\partial n_a}{\partial\bm v_a}\cdot\frac{\partial\bm v_a}{\partial\bm m_a}
=\frac{\partial\bm v}{\partial\bm M}\,\bm\Theta^0_a=-\big(\dot{\bm x}_a-\bm v_a\big) .
\stp
Freeze this vector at the start of the stochastic step, as the gradient energy is frozen now (\texttt{gradient\_freeze},
\texttt{changerep.tex}, Appendix~A), and let the volume follow the momentum linearly in the entropy of the accept/reject:
\st
\boxed{\;
S_a\Big(e_a,\ \bm m_a',\ V_a-\frac{\dot{\bm x}_a-\bm v_a}{p_a}\cdot\big(\bm m_a'-\bm m_a\big)\Big) .\;}
\label{Sfix}
\stp
To first order the entropy change of a move is then
\st
\delta S_a=-\frac{\bm v_a}{T_a}\cdot\delta\bm m_a+\frac{p_a}{T_a}\,\delta V_a=-\frac{\dot{\bm x}_a}{T_a}\cdot\delta\bm m_a ,
\stp
the drift velocity in place of the density-frame velocity: the dissipative bracket uses $\partial H/\partial\bm m_a$, and the work vanishes.
The target $\exp\sum_aS_a$ is a fixed function of the momenta during the sweep, so the Metropolis step is exact for it.  The energy
error at the hand-off is of second order,
\st
\delta H\simeq\tfrac12\sum_ap_a\,\delta\bm m_a\cdot\frac{\partial^2V_a}{\partial\bm m_a\partial\bm m_a}\,\delta\bm m_a ,
\stp
of order $\Delta t$ per step and of the size of the curvature of the kernel stress.  Cost: one frozen vector per particle and one shifted
volume in each entropy evaluation.  The back map is \Eq{Sfix} at the accepted momenta.  The deterministic paths that use
$\beta_a=\texttt{dotx}/\texttt{Teff}$ would take $\dot{\bm x}_a/T_a$ instead.

\paragraph{Files.}
\begin{itemize}
\item \texttt{src/sph.jl}: accept \texttt{ideal = :comoving} with its restrictions.
\item \texttt{src/comoving.jl} (new, loaded by \texttt{src/load.jl}): the density pass, the fixed point, the two loops, the EOS derivatives,
  \texttt{comoving\_leapfrog\_step!}.
\item \texttt{src/ideal.jl}: the branch in \texttt{smooth\_entropy\_energy!} and the dispatch in \texttt{leapfrog\_step!}.
\item \texttt{src/stochastic.jl}, only for \texttt{:comoving}: \texttt{stochastic\_frozen} also freezes the reference momenta and the
  vectors of \Eq{Sfix},
  \st
  \bm m_a^0=\bm m_a,\qquad \bm g_a=-\frac{\dot{\bm x}_a-\bm v_a}{p_a} ,
  \stp
  from the drift loop of the ideal step; every entropy evaluation (\texttt{ask\_star!}, the residual sweep, the back map) takes the volume
  \st
  V_a(\bm m)=V_a+\bm g_a\cdot\big(\bm m-\bm m_a^0\big)
  \stp
  in place of the frozen $V_a$.  The reference is the momentum at the start of the step, not the current one, so the target is one fixed
  function of the momenta through all the sweeps.  For the other options $\bm g_a=0$ and nothing changes.
\item \texttt{test/testcomoving.jl}: the tests below; \texttt{test/run.sh all} with the default bit-identical to \texttt{main}.
\end{itemize}

\paragraph{Tests.}
\begin{enumerate}
\item The EOS derivatives against finite differences of \texttt{solve\_eos!}.
\item $\bm D_{ab}$ and the velocity derivative against finite differences of $N_a$ (done in Python for C2; repeat in Julia for C2, C4, C6).
\item Hamilton's equations: the drift and the force against finite differences of $H$, with the fixed point re-solved, on a jittered
  moving state.
\item The torque identity \Eq{wardN} per star, and $\sum_a\dot{\bm m}_a=0$ to roundoff.
\item $\Delta H$ over a run at $\Delta t$ and $\Delta t/2$; the uniform moving fluid stays uniform.
\item The hand-off error $\delta H$ per stochastic step, with and without \Eq{Sfix}, in an equilibrium ensemble: first order without,
  second order (scaling as $\Delta t$) with.
\end{enumerate}

\paragraph{Cost.}  Each density evaluation becomes two or three passes with an EOS solve each, and the drift needs the second loop's
data, so the ideal step costs roughly twice the present one.  The dissipative steps are unchanged.

%=============================================================================
\section{Checks}
\label{sec:checks}

Ideal step alone, against the present step with the same kernel and $h/a$:
\begin{enumerate}
\item \emph{Derivatives}: \Eq{D} and \Eq{dndv} against finite differences of $N_a$ on jittered stars at $v=0$, $0.3$, $0.8$.
\item \emph{At rest}, uniform fluid on the lattice: identical to the present step.
\item \emph{Uniform moving fluid}, $v_0=0.3$, $0.8$, on the lattice and jittered: the closure \Eq{closure} against $n_a\one$; the size of
  $\bm\Theta^0_a$ and of the drift correction; the state stays uniform.
\item \emph{Conservation}: $\bm P$ to roundoff; $\Delta H$ and $\Delta L$ per step and over the run, against the present step, in a moving
  fluid, a rotating patch, and an equilibrium ensemble.
\item \emph{Sound along and across the flow}: speed and Doppler shift against \texttt{boost.tex}, and against the present step.
\end{enumerate}
With the viscous or stochastic step: the boost benchmark (\texttt{boost.tex}) and the equal-time correlators of a moving fluid, against
\texttt{commonkernel.tex} (\texttt{SPH\_DISSIPATION=common}) and the present scheme.

\end{document}
